\documentclass[10pt]{article}
\usepackage[T1]{fontenc}
\usepackage[utf8]{inputenc}
\usepackage[french]{babel}
\usepackage{newpxtext,newpxmath}
\usepackage[a4paper,top=1.5cm,bottom=1.35cm,left=1.6cm,right=1.6cm]{geometry}
\usepackage{microtype}
\usepackage{xcolor}
\usepackage{titlesec}
\usepackage[most]{tcolorbox}
\usepackage{enumitem}
\usepackage{array}
\usepackage{booktabs}
\usepackage{fancyhdr}

\definecolor{navy}{HTML}{1F3A5F}
\definecolor{grey}{HTML}{54606E}
\definecolor{green}{HTML}{2C6E49}
\definecolor{amber}{HTML}{9C6A16}
\definecolor{red}{HTML}{A6322E}
\definecolor{plum}{HTML}{7A4EA8}
\definecolor{boxgray}{HTML}{F3F5F8}
\definecolor{boxred}{HTML}{FCF2F1}
\definecolor{boxgreen}{HTML}{F1F7F3}
\definecolor{boxplum}{HTML}{F7F3FB}
\definecolor{boxcode}{HTML}{F7F8FA}
\definecolor{boxedge}{HTML}{D9E0EA}

\setlength{\parindent}{0pt}
\setlength{\parskip}{2.6pt}
\linespread{1.0}

\pagestyle{fancy}\fancyhf{}
\lhead{\small\color{grey}Scella}
\rhead{\small\color{grey}Les plus rachetables -- périmètre des deux documents -- 11 août 2026}
\cfoot{\small\color{grey}\thepage}
\renewcommand{\headrulewidth}{0.3pt}

\titleformat{\section}{\color{navy}\large\bfseries}{}{0pt}{}
\titlespacing{\section}{0pt}{11pt}{3pt}
\titleformat{\subsection}{\color{navy}\bfseries}{}{0pt}{}
\titlespacing{\subsection}{0pt}{7pt}{1.5pt}

\newtcolorbox{these}{colback=boxgray,colframe=boxedge,boxrule=0.5pt,arc=2.5pt,
  left=8pt,right=8pt,top=4.5pt,bottom=4.5pt}
\newtcolorbox{alerte}{colback=boxred,colframe=red!55,boxrule=0.8pt,arc=2.5pt,
  left=8pt,right=8pt,top=4.5pt,bottom=4.5pt}
\newtcolorbox{acquis}{colback=boxgreen,colframe=green!45,boxrule=0.8pt,arc=2.5pt,
  left=8pt,right=8pt,top=4.5pt,bottom=4.5pt}
\newtcolorbox{piste}{colback=boxplum,colframe=plum!45,boxrule=0.8pt,arc=2.5pt,
  left=8pt,right=8pt,top=4.5pt,bottom=4.5pt}
\newtcolorbox{extrait}{colback=boxcode,colframe=boxedge,boxrule=0.5pt,arc=1.5pt,
  left=7pt,right=7pt,top=3.5pt,bottom=3.5pt}

% \textcolor plutôt que {\color …} : en tête de cellule p{}, \color insère un whatsit
% qui ouvre une ligne vide et décale verticalement toute la colonne.
\newcommand{\best}[1]{\textcolor{green}{\textbf{#1}}}
\newcommand{\pire}[1]{\textcolor{red}{\textbf{#1}}}
\newcommand{\sz}[1]{\textcolor{grey}{\small\textsf{#1}}}
\newcommand{\code}[1]{\texttt{\small #1}}
\newcommand{\x}{$\times$}
\newcommand{\loc}[1]{{\footnotesize\color{grey}\texttt{#1}}}
\newcommand{\pc}{\,\%}

\begin{document}

\begin{center}
{\Huge\bfseries\color{navy} Scella}\\[4pt]
{\large\color{grey}\itshape Les groupes les plus rachetables --- dans le périmètre des deux documents}\\[4pt]
{\small\color{grey} 11 août 2026 \; · \; \code{Scella\_EN-1.pdf} \x{} \code{Scella\_Produit\_Unique} \; · \; croisement de ce qui est promis et de ce que cela coûte}
\end{center}
\vspace{2pt}

\begin{these}
\textbf{Le résultat central, et il tient à une vérification d'une ligne.} Le seul résultat économique
mesuré en conditions réelles des deux documents --- \emph{« $-$27\pc{} de coût agrégé, jusqu'à
$-$90\pc, latence jusqu'à 3\x{} inférieure, sur Claude + Gemini »} --- est attribué au
\textbf{méta-essaim}. Or Claude et Gemini n'existent que dans \code{efe-swarm-gateway} :
\code{efe-meta-swarm-server} et \code{efe-meta-swarm} en comptent \textbf{zéro occurrence}. Ce chiffre
appartient donc à la \textbf{passerelle} --- c'est-à-dire au binaire qui héberge le \textbf{Relais}.
Corrigée, l'attribution déplace le meilleur argument commercial vers le groupe que le document appelle
déjà lui-même \emph{« la brique la plus finie »}, et que le second PDF a mesuré à \textbf{11 crates
déjà découplés}. Et le périmètre fermé fait apparaître autre chose : \textbf{la pile preuve
réglementaire --- l'actif que j'avais classé premier hier --- est totalement absente de
\code{Scella\_EN}}.
\end{these}

\section{1 · Ce que « dans ces PDF » change au périmètre}

La contrainte est utile plutôt que restrictive, parce que les deux documents ne disent pas la même
chose et ne se recouvrent presque pas. \code{Scella\_EN} dit \textbf{quelle brique résout quel
problème}, avec quelle maturité, et \textbf{ce qui a été mesuré}. \code{Scella\_Produit\_Unique} dit
\textbf{ce que chaque groupe coûte à livrer}, en nomenclature mesurée par \code{cargo tree}. Ni l'un ni
l'autre ne répond seul à la question posée ; croisés, ils y répondent complètement.

Je conserve la discipline des rapports précédents : \textbf{aucun acquéreur n'est nommé, aucun marché
n'est postulé}. C'est d'autant plus facile ici que \code{Scella\_EN} définit ses propres problèmes,
les numérote, et étiquette chaque brique par les numéros qu'elle résout. Je me contente donc de lire
son étiquetage et de le croiser avec les empreintes de l'autre document.

\begin{extrait}
\textbf{Les cinq problèmes, tels que le document les énonce.}
\textbf{(0)} prouver à un tiers, hors ligne, sans vos données ni la parole de l'opérateur ·
\textbf{(1)} orchestrer plusieurs fournisseurs avec une trace auditable ·
\textbf{(2)} garder la donnée chez son propriétaire ·
\textbf{(3)} des décisions rejouables et re-vérifiables ·
\textbf{(4)} réduire coût et hallucination.
\end{extrait}

\section{2 · La carte croisée : brique \x{} problème \x{} maturité \x{} empreinte}

Les colonnes « résout » et « maturité » sont lues dans \code{Scella\_EN} ; la colonne « empreinte »
vient des mesures de \code{Produit\_Unique}. Rien n'est de mon appréciation dans ce tableau.

\renewcommand{\arraystretch}{1.22}
\begin{center}\small
\begin{tabular}{@{}>{\raggedright\arraybackslash}p{3.5cm} c >{\raggedright\arraybackslash}p{7.0cm} >{\raggedright\arraybackslash}p{3.7cm}@{}}
\toprule
\textbf{Brique} \sz{(Scella\_EN)} & \textbf{Résout} & \textbf{Maturité --- dite par le document} & \textbf{Empreinte} \sz{(Produit\_Unique)} \\
\midrule
\textbf{Le Relais inter-modèles} & (1) & \best{« la brique la plus finie »} --- codée et poussée de bout en bout, servie opt-in, $\sim$220 tests, 4 revues adverses, B.0$\rightarrow$B.4 & \best{11 crates}, déjà découplé \\
\textbf{Le méta-essaim} & (1)(2) & « réel et testé », mTLS strict, manifeste moteur porté par le protocole & 7 défaut / \pire{26 réel} / $\sim$9 corrigé \\
\textbf{Le reçu} & (0) & la fondation --- les six briques n'écrivent qu'un seul reçu & présent dans toutes \\
\textbf{Le Cadrage} & (4) & \pire{« architecture écrite, prête à coder »} & \pire{0 ligne} \\
\textbf{L'essaim médiateur} & (1)(3) & « deux tiers déjà là » ; la chaîne déclarée n'est pas codée & \pire{0 ligne} \\
Nano-robots $+$ $\Phi$ & (3) & mesuré, déterministe, $\hat{L}$ empirique & \code{efe-core} $=$ \best{1 crate} \\
Flexible MLP & (4) & « pièce séparée, non intégrée » & hors dépôt \\
\bottomrule
\end{tabular}
\end{center}

\section{3 · Le défaut d'attribution qui déplace le verdict}

C'est la vérification la plus importante de ce rapport, et elle suffit à réordonner le classement.

\code{Scella\_EN} place son \textbf{seul résultat économique mesuré en conditions réelles} sur le
méta-essaim :

\begin{extrait}
\sz{« \textbf{Meta-swarm (coordination)} --- validated live on real providers (\textbf{Claude + Gemini}) :
cost cut ($\sim$$-$27\pc{} aggregate, up to $\sim$90\pc{} on the cheap-routable), quality held, latency up
to $\sim$3\x{} lower. »}
\end{extrait}

Or, mesuré dans le code : \textbf{Claude et Gemini n'existent que dans \code{efe-swarm-gateway}} ---
six fichiers. Et \code{efe-meta-swarm-server} comme \code{efe-meta-swarm} comptent \textbf{zéro
occurrence} de Claude, Anthropic, OpenAI ou Gemini.

\begin{alerte}
\textbf{Le méta-essaim de la Section I ne connaît aucun fournisseur frontière.} Il est
Tensor$\rightarrow$Tensor sur gRPC, et c'est précisément ce que le document en dit par ailleurs. Le
$-$27\pc{} \textbf{n'a pas pu être produit par lui} : il appartient à la \textbf{passerelle}, c'est-à-dire
au binaire qui héberge le \textbf{Relais}.

Ce n'est pas une scorie rédactionnelle. C'est le \emph{seul} chiffre d'argent des deux documents, et il
est adossé à un composant qu'un examinateur disqualifie en un \code{grep}. Corrigée, l'attribution
\textbf{transfère le meilleur argument commercial} du groupe B vers le groupe A.
\end{alerte}

Il faut être juste sur l'origine de la confusion : le dépôt contient bien deux couches de coordination
distinctes --- la frontière LLM dans la passerelle, et le méta-essaim latent en gRPC --- et le document
les décrit toutes deux honnêtement, séparément. Ce qui manque n'est pas l'honnêteté, c'est le
rattachement du résultat à la bonne couche.

\section{4 · Le classement}

Quatre groupes constructibles à partir de ces deux documents ; rangs de 1 à 4 par critère.

\renewcommand{\arraystretch}{1.2}
\begin{center}\small
\begin{tabular}{@{}r >{\raggedright\arraybackslash}p{5.5cm} c ccccc c@{}}
\toprule
& \textbf{Le groupe} & \sz{Résout} & \sz{Chiffré} & \sz{Matur.} & \sz{Empr.} & \sz{Sans croy.} & \sz{Doc/code} & \textbf{Moy} \\
\midrule
\textbf{1} & \textbf{A · Le relais souverain} \sz{Relais + frontière + reçu} & \textbf{(0)(1)(2)} & \best{1} & \best{1} & \best{1} & \best{1} & 2 & \best{1,20} \\
2 & C · Le Cadrage + reçu & (0)(4) & 2 & \pire{4} & \pire{4} & 2 & \best{1} & \textbf{2,60} \\
3 & B · L'essaim attestable \sz{méta-essaim + reçu} & (0)(1)(2) & 3 & 2 & 3 & 3 & 3 & \textbf{2,80} \\
4 & D · La lignée latente \sz{nano + $\Phi$ + méta} & (2)(3) & \pire{4} & 3 & 2 & \pire{4} & \pire{4} & \pire{3,40} \\
\bottomrule
\end{tabular}
\end{center}

\textbf{A gagne sur quatre critères sur cinq}, et chacune des raisons est écrite dans l'un des deux
documents plutôt que dans mon appréciation.

\textbf{C'est le document lui-même qui l'appelle « la brique la plus finie »} --- codée et poussée de
bout en bout, servie opt-in, $\sim$220 tests, quatre revues adverses, vingt-six correctifs. Aucune autre
brique ne reçoit ce qualificatif, et c'est le genre de phrase qu'un auteur n'écrit pas à la légère sur
son propre travail.

\textbf{C'est là que vit l'argent}, une fois l'attribution du §3 corrigée : $-$27\pc{} agrégé, jusqu'à
$-$90\pc{} sur le routable-bon-marché, latence 3\x{} inférieure, qualité tenue --- sur de vrais
fournisseurs, pas en simulation.

\textbf{\code{Produit\_Unique} l'a mesuré à 11 crates, déjà découplé, sans monolithe.} C'est le seul
candidat qui n'exige \textbf{aucun correctif de packaging}, là où B traîne ses 26 crates tant que la
recette Mode-1 n'est pas extraite d'\code{efe-server}.

\textbf{Il n'exige aucune croyance.} « Faire coopérer Claude, GPT et un modèle interne sans colle
tierce, avec une trace auditable » s'évalue sans adhérer à quoi que ce soit sur les essaims latents ---
et un pari est exactement ce qu'une revue d'acquisition cherche à éliminer.

\begin{alerte}
\textbf{C arrive deuxième, et ce rang est trompeur --- il faut le dire.} Le Cadrage a un résultat
mesuré ($-$85\pc{} de tokens d'entrée, 132\,$\rightarrow$\,20, réponse vérifiée SUPPORTED), il ne demande
aucune croyance, et le document est \emph{franc} sur son état, ce qui lui vaut le premier rang en
cohérence. Mais votre question dit \textbf{« maintenant »}, et le Cadrage compte \textbf{zéro ligne}.
Son rang mesure une \emph{promesse}, pas une \emph{disponibilité}. Il ne peut pas être racheté
aujourd'hui, quelle que soit sa moyenne.
\end{alerte}

\textbf{Et B tombe troisième}, alors que mes deux rapports précédents le couronnaient. Dans le
périmètre de ces documents, son argument distinctif --- le $-$27\pc{} --- \textbf{ne lui appartient
pas}, et ce qui reste exige la thèse latente. C'est la troisième fois de la série qu'un changement
d'axe déplace le vainqueur ; c'est le signe que les axes mesurent réellement des choses différentes.

\section{5 · Le résultat principal : ce que les deux documents \emph{ne contiennent pas}}

C'est la conclusion la plus utile, et elle ne pouvait apparaître qu'en fermant le périmètre comme vous
l'avez fait. J'ai cherché dans \code{Scella\_EN} les termes de l'actif que j'avais classé
\textbf{premier en acquérabilité} la veille.

\renewcommand{\arraystretch}{1.15}
\begin{center}\small
\begin{tabular}{@{}>{\raggedright\arraybackslash}p{9.8cm} c@{}}
\toprule
\textbf{Terme cherché dans \code{Scella\_EN}} & \textbf{Occurrences} \\
\midrule
\code{reg-evidence} · \code{audit-api} · \code{provenance-capture} · \code{forensic} & \pire{0} \\
\textbf{AI Act} · \textbf{GDPR} · \textbf{CSDDD} & \pire{0} \\
« zero-dependency » · « tap » & \pire{0} \\
\bottomrule
\end{tabular}
\end{center}

\begin{acquis}
\textbf{La pile preuve réglementaire --- huit crates, 12 862 lignes, 240 tests, trois obligations
légales nommées --- est totalement absente du document qui présente Scella.}
\end{acquis}

Un acquéreur qui ne lit que ces deux PDF \textbf{n'apprendra jamais qu'elle existe}. Ce n'est donc pas
un défaut de classement de ma part : dans le périmètre imposé, elle n'est pas classable, et je ne l'ai
pas classée. Mais son absence \emph{est} le résultat --- le document de présentation omet la pièce qui
répond à trois obligations non optionnelles, au profit de six briques dont \textbf{deux comptent zéro
ligne}. C'est le déséquilibre le plus coûteux des deux documents, et il ne se corrige pas par du code.

\section{6 · Deux péremptions qu'un examinateur trouvera seul}

\textbf{« external anchor not yet wired »} \sz{(p.\,3)}. L'ancre \emph{interne} \textbf{a été câblée} ---
journal append-only, troncature détectable \emph{sans clé}. Ce qui manque est l'ancrage \textbf{tiers}.
La formulation juste serait : \emph{« internal anchor wired ; third-party anchor not »}. En l'état, le
document \textbf{sous-vend} une levée réelle, ce qui est le défaut symétrique de la survente et coûte
tout aussi cher.

\textbf{« nearly 1,200 green tests »} \sz{(p.\,4)}. Portée non définie --- sur quels crates ? Je n'ai
jamais pu la re-mesurer, et je l'avais déjà signalé. Un chiffre non reproductible dans un document par
ailleurs rigoureux est exactement ce qui appelle la question suivante, et la question suivante porte
rarement sur le chiffre.

\section{7 · Le jugement}

\begin{acquis}
\textbf{Dans le périmètre de ces deux documents, le groupe le plus rachetable maintenant est A --- le
relais souverain : le Relais inter-modèles, la coordination frontière, et le reçu.}
\end{acquis}

Onze crates déjà découplés ; la brique que le document lui-même désigne comme la plus finie ; le seul
résultat économique mesuré sur de vrais fournisseurs ; et une promesse qui n'exige d'adhérer à aucune
thèse. Il couvre \textbf{(0)}, \textbf{(1)} et une part de \textbf{(2)} --- soit trois des cinq
problèmes que le document s'est lui-même donnés.

\textbf{Trois actions, par ordre de rendement, et aucune n'est du code :}

\begin{enumerate}[leftmargin=1.6em,itemsep=2.5pt,topsep=3pt]
\item \textbf{Corriger l'attribution du $-$27\pc} --- le déplacer du méta-essaim vers la couche qui l'a
      produit. C'est le chiffre le plus fort des deux documents, et il est aujourd'hui \emph{invérifiable
      là où il est écrit}.
\item \textbf{Ajouter la pile preuve au document.} Trois obligations légales nommées, un producteur
      générique, un vérifieur sans dépendance : hors périmètre, c'est l'actif que j'ai classé premier
      --- et il n'apparaît nulle part.
\item \textbf{Corriger les deux péremptions} --- l'ancre \emph{est} câblée, et le chiffre de tests doit
      être mesurable ou retiré.
\end{enumerate}

\begin{these}
\small\textbf{Rappel de cadrage, inchangé dans toute la série.} Le reçu atteste l'\textbf{intégrité}
d'une suite de pas, jamais la \textbf{correction} d'une décision. Il est \emph{tamper-evident pour un
détenteur de clé}, pas \emph{trustless pour un tiers} --- ce que \code{Scella\_EN} dit d'ailleurs
correctement. Les participants du mode réel exécutent une fusion latente Mode-1 numérique : ce ne sont
pas des « vrais SLM ». Et ce document évalue des \emph{propriétés d'acquérabilité} : il ne prédit aucun
acquéreur et n'en nomme aucun.
\end{these}

\section{Annexe · Vérifier chaque affirmation soi-même}

Les \emph{rangs} sont un jugement. Les \textbf{citations} et les \textbf{zéros} se rejouent --- et ce
sont eux qui portent le verdict.

\renewcommand{\arraystretch}{1.22}
\begin{center}\small
\begin{tabular}{@{}>{\raggedright\arraybackslash}p{5.5cm} >{\raggedright\arraybackslash}p{10.6cm}@{}}
\toprule
\textbf{Affirmation} & \textbf{Ce qui la rejoue} \\
\midrule
Le texte des deux documents &
  \code{pdftotext -layout Scella\_EN-1.pdf -} \\
Claude/Gemini ne vivent que dans la passerelle &
  \code{grep -rln 'Gemini' crates -{}-include='*.rs'} \\
Le méta-essaim ignore les fournisseurs &
  \code{grep -rlin 'claude\textbackslash|anthropic\textbackslash|openai\textbackslash|gemini' \textbackslash} \newline
  \code{\ \ crates/efe-meta-swarm-server/src crates/efe-meta-swarm/src} \sz{$\rightarrow$ vide} \\
La pile preuve est absente du document &
  \code{pdftotext -layout Scella\_EN-1.pdf - | grep -ci 'AI Act\textbackslash|GDPR\textbackslash|CSDDD'} \sz{$\rightarrow$ 0} \\
Les empreintes 11 / 7 / 26 crates &
  \code{Scella\_Produit\_Unique\_2026-08-11.pdf}, §2 --- et son annexe donne la commande \\
Les deux péremptions &
  \code{pdftotext -layout Scella\_EN-1.pdf - | grep -n 'anchor not yet\textbackslash|1,200'} \\
\bottomrule
\end{tabular}
\end{center}

Un avertissement pour finir. Ce rapport est \textbf{borné aux deux documents fournis} : il ne classe pas
le dépôt, et un actif absent de \code{Scella\_EN} y est absent par construction, non par jugement. Le
fichier analysé est \code{Scella\_EN-1.pdf} (4 pages, 59 901 octets, 14 août 01\,h\,54) --- il diffère de
\code{Scella\_EN.pdf} du 10 août, et les deux péremptions du §6 lui appartiennent en propre.

\end{document}
